\textbf{Keywords:} TempSense, Node Voltage, Ground, VDD, Clocks,
Digital, Bias, RESET (POR), Package, Why ESD, CDM (Gauss, INV), HBM
(01,10,02,20,12,21**), GGNMOS, Latch-up

Video is from 2024, so the plan might not be exactly the same. In
addition, we're not using Caravel for the tapeout, but rather
TinyTapeout.

\section{What blocks must our IC
include?}\label{what-blocks-must-our-ic-include}

The project is to design an integrated temperature sensor.

First, we need to have an idea of what comes in and out of the
temperature sensor. Before we have made the temperature sensor, we need
to think what the signal interface could be, and we need to learn.

Maybe we read
\href{http://ei.ewi.tudelft.nl/docs/TSensor_survey.xls}{Kofi Makinwa's
overview of temperature sensors} and find one of the latest papers,

A BJT-based CMOS Temperature Sensor with Duty-cycle-modulated Output and
±0.54 °C (3-sigma) Inaccuracy from -40 °C to 125 °C
\citeproc{ref-huang21}{{[}1{]}}.

At this point, you may struggle to understand the details of the paper,
but at least it should be possible to see what comes in and out of the
module. What I could find is in the table below, maybe you can find
more?

\begin{longtable}[]{@{}lllll@{}}
\toprule\noalign{}
Pin & Function & in/out & Value & Unit \\
\midrule\noalign{}
\endhead
\bottomrule\noalign{}
\endlastfoot
VDD\_3V3 & analog supply & in & 3.0 & V \\
VDD\_1V2 & digital supply & in & 1.2 & V \\
VSS & ground & in & 0 & V \\
CLK\_1V2 & clock & in & 20 & MHz \\
RST\_1V2 & digital & out & 0 or 1.2 & V \\
I\_C & bias & in & ? & uA? \\
PHI1\_1V2 & digital & out & 0 or 1.2 & V \\
PHI2\_1V2 & digital & out & 0 or 1.2 & V \\
DCM\_1V2 & digital & out & 0 or 1.2 & V \\
\end{longtable}

This list contains supplies, clocks, digital outputs, bias currents and
a ground. Let me explain what they are.

\subsection{Supply}\label{supply}\index{supply@Supply}

The temperature sensor has two supplies, one analog (3.3 V) and one
digital (1.2 V), which must come from somewhere.

We're using \href{http://www.tinytapeout.com}{TinyTapeout}

That has ability for both 3.3 V and 1.8 V. An external low dropout
regulator (LDO) provides the digital supply (1.8 V).

See for example
\href{https://caravel-harness.readthedocs.io/en/latest/maximum-ratings.html}{absolute
maximum ratings}

\subsubsection{Ground}\label{ground}

Most ICs have a ground, a pin which is considered 0 V. It may have
multiple grounds. Remember that a voltage is only defined between two
points, so it's actually not true to talk about a voltage in a node (or
on a wire). A voltage is always a differential to something. We've (as
in global electronics engineers) have just agreed that it's useful to
have a ``node'' or ``wire'' we consider 0 V.

\subsubsection{Clocks}\label{clocks}

Most digital need a clock, and TinyTapeout can provide a 50 MHz clock
which should suffice for most things. We could probably just use that
clock for our temperature sensor.

\subsubsection{Digital}\label{digital}

We need to read the digital outputs. We could either feed those off
chip, or use a on chip micro-controller. The TinyTapeout includes
options to do both. We could connect digital outputs to the logic
analyzer, and program the MCU to store the readings. Or we could connect
the digital output to the I/O and use an instrument in the lab.

\subsubsection{Bias}\label{bias}

The TinyTapeout does not provide bias currents (that I found), so that
is something you will need to make.

\subsubsection{Conclusion}\label{conclusion}

Even a temperature sensor needs something else on the IC. We need
digital input/output, clock generation (PLL, oscillators), bias current
generators, and voltage regulators (which require a constant reference
voltage).

I would claim that any System-On-Chip will always need these blocks!

I want you to pause, take a look at the

\href{https://analogicus.com/aic2026/plan/}{course plan}

and now you might understand why I've selected the topics.

\subsubsection{One more thing}\label{one-more-thing}

There is one more function we need when we have digital logic and a
power supply. We need a ``RESET'' system.

Digital logic has a fundamental assumption that we can separate between
a ``1'' and a ``0'', which is usually translated to for example 1.8 V
(logic 1) and 0 V (logic 0). But if the power supply is at 0 V, before
we connect the battery, then that fundamental assumption breaks.

When we connect the battery, how do we know the fundamental assumption
is OK? It's certainly not OK at 30 mV supply. How about 500 mV? or 1.0
V? How would we know?

Most ICs will have a special analog block that can keep the digital
logic, bias generators, clock generators, input/output and voltage
regulators in a \textbf{safe} state until the power supply is high
enough (for example 1.62 V).

One of the challenges with a Power On Reset (POR) is that we want to
keep the system in a reset state until we're sure that the power is on.
Another challenge is that the POR should not consume current.

If we make a level triggered (triggers when VDD reaches a certain
level), then we need a reference, a comparator and maybe other circuits.
As a result, potentially high current.

If we make a delay based POR, then we need a long delay, which means
large resistors or capacitors. Accordingly, high cost.

Below is an idea for a
\href{https://patents.google.com/patent/GB2509147A/en?inventor=carsten+wulff&oq=carsten+wulff}{Power-On-Reset}
I had way back when. The POR uses a delay based on the tunneling current
in a thin oxide transistor (2), and uses a thick-oxide transistor (3) as
a capacitor. The output X would go to a Schmitt trigger (5).


{
\centering
\aicfig{media/por.png}
}

\emph{Figure 1: Power On Reset using gate tunneling }

\section{Electrostatic Discharge}\label{electrostatic-discharge}\index{electrostatic discharge@Electrostatic discharge}

If you make an IC, you must consider Electrostatic Discharge (ESD)
Protection circuits

ESD events are tricky. They are short (ns), high current (Amps) and
poorly modeled in the SPICE model.

Most SPICE models will not model correctly what happens to an transistor
during an ESD event. The SPICE models are not made to model what happens
during an ESD event, they are made to model how the transistors behave
at low fields and lower current.

But ESD design is a must, you have to think about ESD, otherwise your IC
will never work.

Consider a certain ESD specification, for example 1 kV human body model,
a requirement for an integrated circuit.

By requirement I mean if the 1 kV is not met, then the project will be
delayed until it is fixed. If it's not fixed, then the project will be
infinitely delayed, or in other words, canceled.

Now imagine it's your responsibility to ensure it meets the 1 kV
specification, what would you do? I would recommend you read one of the
few ESD books in existence - see the reading list at the end of the
chapter - and rely on your understanding of PN-junctions.

The industry has agreed on some common test criteria for electrostatic
discharge. Test that model what happens when a person touches your IC,
during soldering, and PCB mounting. If your IC passes the test then it's
probably going to survive in volume production

Standards for testing at
\href{https://www.jedec.org/category/technology-focus-area/esd-electrostatic-discharge-0}{JEDEC}

The JEDEC standard splits ESD events into Human body model, Charged
device model, and System level ESD.

Once mounted on the PCB, the ICs can be more protected against ESD
events, however, it depends on the PCB, and how that reacts to a
current.

Take a look at your USB-A connector, you will notice that the outer
pins, the power and ground, are made such that they connect first, The
\(D+\) and \(D-\) pins are a bit shorter, so they connect some \(\mu\)s
later. The reason is ESD. The power and ground usually have a low
impedance connection in decoupling capacitors and power circuits, so
those can handle a large ESD zap. The signals can go directly to an IC,
and thus be more sensitive.

We won't go into details on System level ESD, as that is more a PCB type
of concern. The physics are the same, but the details are different.

\subsection{Human body model (HBM)}\label{human-body-model-hbm}\index{human body model (hbm)@Human body model (HBM)}

HBM is the ``simple'' version of ESD, a model can be seen in Figure 2.
Some of the properties of HBM are:

\begin{itemize}
\tightlist
\item
  Models a person touching a device with a finger
\item
  \textbf{Long} duration (around 100 ns)
\item
  Acts like a current source into a pin
\item
  Can usually be handled in the I/O ring
\item
  4 kV HBM ESD is 2.67 A peak current
\end{itemize}


{
\centering
\aicfig{media/esd_hbm_finger.png}
}

\emph{Figure 2: Human body model (HBM)}

More on circuits that protect from HBM later.

\subsection{Charged device model (CDM)}\label{charged-device-model-cdm}\index{charged device model (cdm)@Charged device model (CDM)}

\begin{quote}
An IC left alone for long enough will equalize the Fermi potential
across the whole IC.
\end{quote}

Not entirely a true statement, but roughly true. One exception is
non-volatile memory, like flash, which uses
\href{https://en.wikipedia.org/w/index.php?title=Field_electron_emission&oldformat=true\#Fowler–Nordheim_tunneling}{Fowler-Nordheim}
tunneling to charge and discharge a capacitor that keeps its charge for
a very, very long time.

I'm pretty sure that if you leave an SSD hardrive to the
\href{https://en.wikipedia.org/wiki/Heat_death_of_the_universe}{heat
death of the universe} in maybe \(10^{10^{10^{56}}}\) years, then the
charges will equalize, and the Fermi level will be the same across the
whole IC, so it's just a matter of time.

Assume there is an equal number of electrons and protons on the IC.
According to Gauss' law

\[ \oint_{\partial \Omega} \mathbf{E} \cdot d\mathbf{S} = \frac{1}{\epsilon_0} \iiint_{V} \rho
\cdot dV\]

Which says that the electric field through the surface is the volume
integral of the charges inside the surface. If there are the same amount
of protons and electrons, and the distribution is even, then there will
be no field through IC surface. As such, there is no external electric
field from the IC.

If we place an IC in an electric field, the charges inside will
redistribute. Flip the IC on its back, place it on an metal plate with
an insulator in-between, and charge the metal plate to 1 kV, as shown in
Figure 3.


{
\centering
\aicfig{media/cdm.png}
}

\emph{Figure 3: Charged Device Model (CDM) testing}

Inside the integrated circuit, electrons and holes will redistribute to
compensate for the electric field. Closest to the metal plate there will
be a negative charge, and furthest away there will be a positive charge.

This comes from the fact that if you leave a metal inside an electric
field for long enough the metal will not have any internal field. If
there was an internal field, the charges would move. Over time the
charges will be located at the ends of the metal.

Take a grounded wire, touch one of the pins on the IC. Since we now have
a metal connection between a pin and a low potential the charges inside
the IC will redistribute extremely quickly, on the order of a few ns.

During this Charged Device Model event the internal fields in the IC
will be chaotic, but at any given point in time, the voltage across
sensitive devices must remain below where the device physically breaks.

Take the MOSFET transistor. Between the gate and the source there is an
thin oxide, maybe a few nm. If the field strength between gate and
source is high enough, then the force felt by the electrons in co-valent
bonds will be \(\vec{F} = q\vec{E}\). At some point the co-valent bonds
might break, and the oxide could be permanently damaged. Think of a
lightning bolt through the oxide, it's a similar process.

Our job, as electronics engineers, is to ensure we put in additional
circuits to prevent the fields during a CDM event from causing damage.

For example, let's say I have two inverters powered by different supply,
VDD1 and VDD2. If I in my ESD test ground VDD1, and not VDD2, I will
quickly bring VDD1 to zero, while VDD2 might react slower, and stay
closer to 1 kV. The gate source of the PMOS in the second inverter will
see approximately 1 kV across the oxide, and will break. How could I
prevent that?


{
\centering
\aicfig{media/esd_cdm_domains_tikz.png}
}

\emph{Figure 4: Cross domain voltage problem with CDM (or indeed HBM)
events }

Assuming some luck, then VDD1 and VDD2 are separate, but the same
voltage, or at least close enough, I can take two diodes, connected in
opposite directions, between VDD1 and VDD2. As such, when VDD1 is
grounded, VDD2 will follow but maybe be 0.6 V higher. As a result, the
PMOS gate never sees more than approximately 0.6 V across the gate
oxide, and everyone is happy.

Now imagine an IC with hundreds of supplies, and billions of inverters.
How can I make sure that everything is OK?

CDM is tricky, because there are so many details, and it's easy to miss
one that makes your circuit break.

\section{An HBM ESD zap example}\label{an-hbm-esd-zap-example}\index{an hbm esd zap example@An HBM ESD zap example}

Imagine a ESD zap between VSS and VDD. How can we protect the device?

The positive current enters the VSS, and leaves via the VDD, so our
supplies are flipped up-side down. It's a fair assumption that none of
the circuits inside will work as intended.

But the IC must not die, so we have to lead the current to ground
somehow.


{
\centering
\aicfig{media/esd_hbm_model.png}
}

\emph{Figure 5: ESD HBM zap example }

Let's simplify and think of the possible permutations, shown in Figure
6. We don't know where the current will enter nor where it will leave
our circuit, so we must make sure that all combinations are covered.


{
\centering
\aicfig{media/esd_perm_tikz.png}
}

\emph{Figure 6: The six ways a zap can cross a three terminal chip.
Three pads means six ordered pairs, and ESD protection is the promise
that every one of them has somewhere for the current to go. The chip is
empty here on purpose; the next three figures fill it in.}

When the current enters VSS and must leave via VDD, then it's simple, we
can use a diode, as shown in Figure 7.

Under normal operation the diode will be reverse biased, and although it
will add some leakage, it will not affect the normal operation of our
IC.


{
\centering
\aicfig{media/esd_zap01_tikz.png}
}

\emph{Figure 7: Protection for a zap from ground to VDD. One diode,
reverse biased whenever the chip is running, forward biased the moment
VSS goes above VDD. The column crosses the pin rail without a junction
dot, which is a statement that nothing is connected there yet.}

The same is true for current in on VSS and out on PIN. Here we can also
use a diode, as shown in Figure 8.


{
\centering
\aicfig{media/esd_zap02_tikz.png}
}

\emph{Figure 8: Protection for a zap from ground to a pin. Nothing new
happens, and that is the point: a signal pin sits above VSS in normal
operation for the same reason VDD does, so the same reverse biased diode
covers it.}

For a current in on VDD and out on VSS we have a challenge. That's the
normal way for current to flow.

For those from Norway that have played a kids game
\href{https://www.youtube.com/watch?v=jtZ1R9_Lu-4}{Bjørnen sover},
that's a apt mental image. We want a circuit that most of the time
sleeps, and does not affect our normal IC operation. But if a huge
current comes in on VDD, and the VDD voltage shoots up fast, the circuit
must wake up and bring the voltage down.

If the circuit triggers under normal operating condition, when you're
watching a video on your phone, your battery will drain very fast, and
your phone might even catch fire.

As such, ESD design engineers have a ``ESD design window''. Never let
the ESD circuit trigger when VDD \textless{} normal, but always trigger
the ESD circuit before VDD \(>\) breakdown of circuit.

A circuit that can sometimes be used, if the ESD design window is not
too small, is the Grounded-Gate-NMOS in Figure 9.


{
\centering
\aicfig{media/esd_rails_all_tikz.png}
}

\emph{Figure 9: All six permutations covered by four devices. Each
device is labelled with the zap it carries, and two of them carry one
jointly: a strike from VDD to the pin goes down the left grounded gate
NMOS to VSS and back up through the pin diode, which is why
\(1\rightarrow2\) appears twice. That is the real structure of an ESD
network. It is not one path per permutation but a small set of devices
arranged so every permutation has a series combination available, which
is why adding a pin costs two diodes rather than six.}

\section{The grounded gate NMOS}\label{the-grounded-gate-nmos}\index{the grounded gate nmos@The grounded gate NMOS}

If you try the circuit in Figure 10 with the normal BSIM spice model, it
will not work. The transistor model does not include that part of the
physics.

We need to think about how electrons, holes PN-junctions and bipolars
work. Let's refresh quantum mechanics a bit.


{
\centering
\aicfig{media/esd_ggnmos_tikz.png}
}

\emph{Figure 10: The grounded gate NMOS (GGNMOS) }

Electrons sticking to atoms (bound electrons), can only exist at
discrete energy levels. As we bring atoms closer to each-other the
discrete energy levels will split, as computed from Schrodinger, into
bands of allowed energy states. These bands of energy can have lower
energy than the discrete energy levels of the atom. That's why some
atoms stick together and form molecules through co-valent bonds, ionic
bonds, or whatever the chemists like to call it. It's all the same
thing, it's lower energy states that make the electrons happy, some are
strong, some are weak.

For silicon the
\href{https://www.iue.tuwien.ac.at/phd/wessner/node31.html}{energy band
structure} is tricky to compute, so we simplify to band diagrams that
only show the lowest energy conduction band and highest energy valence
band.

Electrons can move freely in the conduction band (until they hit
something, or scatter), and electrons moving in the valence band act
like positive particles, nicknamed holes.

How many free charges there are in a band is given by Fermi-Dirac
distribution and the density of states (allowed energy levels).

If an electron, or a hole have sufficient energy (accelerated by a
field), they can free an electron/hole pair when they scatter off an
atom. If you break too many bonds between atoms, your material will be
damaged.

Assume a transistor like the one in Figure 11. The gate, source and bulk
is connected to ground. The drain is connected to a high voltage.


{
\centering
\aicfig{media/esd_ggnmos_xsec_tikz.png}
}

\emph{Figure 11: Cross section of the grounded gate NMOS }

The process of a GGNMOS will be (1) Avalanche, (2) hole accumulation,
(3) forward bias of PN-junction, and (4) direct electron current from
source to drain.

The first thing that can happen is that the field in the depletion zone
between drain and bulk (1) is large, due to the high voltage on drain,
and the thin depletion region.

In the substrate (P-) there are mostly holes, but there are also
electrons. If an electron diffuses close to the drain region it will be
swept across to drain by the high field.

The high field might accelerate the electron to such an energy that it
can, when it scatters off the atoms in the depletion zone, knock out an
electron/hole pair.

The hole will go to the substrate (2), while the new electron will
continue towards drain. The new electron can also knock out a new
electron/hole pair (energy level is set by impact ionization of the
atom), so can the old one assuming it accelerates enough.

One electron turn into two, two to four, four to eight and so on. The
number of electrons can quickly become large, and we have an avalanche
condition. Same as a snow avalanche, where everything was quiet and
nice, now suddenly, there is a big trouble.

Usually the avalanche process does not damage anything, at least
initially, but it does increase the hole concentration in the bulk. The
number of holes in the bulk will be the same as the number of electrons
freed in the depletion region.

The extra holes underneath the transistor will increase the local
potential. If the substrate contact (5) is far away, then the local
potential close to the source/bulk PN-junction (3) might increase enough
to significantly increase the number of electrons injected from source.

Some of the electrons will find a hole, and settle down, while others
will diffuse around. If some of the electrons gets close to the drain
region, and the field in the depletion zone, they will be accelerated by
the drain/bulk field, and can further increase the avalanche condition.

For a normal transistor, not designed to survive, the electron flow (4)
can cause local damage to the drain. Normally there is nothing that
prevents the current from increasing, and the transistor will eventually
die.

If we add a resistor to the drain region (unsilicided drain), however,
we will slow down the electron flow, and we can get a stable condition,
and design a transistor that survives.

Turns out, that every single NMOS has a sleeping bear. A parasitic
bipolar. That's exactly what this GGNMOS is, a bipolar transistor,
although a pretty bad one, that is designed to trigger when avalanche
condition sets in and is designed to survive.

A normal NMOS, however, can also trigger, and if you have not thought
about limiting the electron current, it can die, with IC killing
consequences. Specifically, the drain and source will be shorted by
likely the silicide on top of the drain, and instead of a transistor
with high output impedance, we'll have a drain source connection with a
few kOhm output impedance.

Take a look at New Ballasting Layout Schemes to Improve ESD Robustness
of I/O Buffers in Fully Silicided CMOS Process
\citeproc{ref-ker09}{{[}2{]}} for the pretty pictures you'll get when
the drain/source breaks.

\section{But I just want a digital input, what do I
need?}\label{but-i-just-want-a-digital-input-what-do-i-need}

Even if it's only a digital input, you still need to consider ESD
events.

Below is a complete digital input network.

Assuming we have a
\href{https://en.wikipedia.org/wiki/Flat_no-leads_package}{QFN package}
package there will be a bond-wire from the package metal, to our die
pad.

Right after the die pad, sometimes under, there will be a primary ESD
protection that can conduct, in all directions, between input, supply
and ground.

From the input it's common to have a resistor to reduce the probability
of currents going towards the core area.

Before we get to a transistor gate oxide it's common to have a set of
secondary protection circuits. A resistor further reduces the current,
and two local clamps (GGPMOS and GGNMOS) ensure that the voltage across
the transistor gate does not go to breakdown levels.


{
\centering
\aicfig{media/esd_input_prot_tikz.png}
}

\emph{Figure 12: Full protection of an input including secondary
protection }

\subsection{Input buffer}\label{input-buffer}\index{input buffer@Input buffer}

An input buffer can be seen below. I like to include a RC low-pass
filter to filter out the RF frequencies (I don't want my input to toggle
if a phone is on top of my circuit).

After the RC filter we need a
\href{https://en.wikipedia.org/wiki/Schmitt_trigger}{Schmitt trigger},
you can find a Schmitt trigger at
\href{https://analogicus.github.io/jnw_tr_sky130A/schematic.html}{JNW\_TR\_SKY130A}.

The Schmitt trigger must be with thick oxide gates and with IO supply
(for example 3.0 V).

The first inverter must also be a thick oxide inverter, however, the
supply of the inverter will be core supply (for example 1.2 V). The
thick oxide inverter provides a level-shift to core supply.

The last inverter is just to get the polarity of the TO\_CORE signal the
same as the input.


{
\centering
\aicfig{media/fig_methodology.png}
}

\emph{Figure 13: Full digital input including RC filter, Schmitt trigger
and level shifters: (a) schematic, (b) netlist, (c) the compiled layout,
and (d) the ciccreator recipe that generated it}

\section{Latch-up}\label{latch-up}\index{latch-up@Latch-up}

Another fun physics problem can happen in digital logic that is close to
an electron source, like a connection to the real world, what we call a
pad. A pad is where you connect the bond-wire in a QFN type of package
with \href{https://en.wikipedia.org/wiki/Wire_bonding}{wire-bonding}

Assume we have the circuit in Figure 14. Under certain conditions we can
get a short from VDD to ground.


{
\centering
\aicfig{media/esd_latchup_tikz.png}
}

\emph{Figure 14: Inverter that suddenly shorts from VDD to ground
located close to a PAD}

Consider the cross section of the inverter in Figure 15. The latch-up
process starts with electron injection (1), then forward bias of PMOS
source/drain junction (2), forward bias of NMOS source/drain junction
(3) , and finally positive feedback .


{
\centering
\aicfig{media/esd_scr_xsec_tikz.png}
}

\emph{Figure 15: Cross section of an inverter }

\subsection{Electron injection}\label{electron-injection}\index{electron injection@Electron injection}

Assume that we have an electron source, for example a pad that is below
ground for a bit. This will inject electrons into the substrate/bulk (1)
and electrons will diffuse around.

If some of the electrons comes close to the N-well depletion region (2)
they will be swept across by the built-in field. As a result, the
potential of the N-well will decrease, and we can forward bias the
source or drain junction of a PMOS.

\subsubsection{Forward biased PMOS source or drain
junction}\label{forward-biased-pmos-source-or-drain-junction}

With a forward biased source/bulk junction (2), holes will be injected
into the N-Well, but similarly to the GGNMOS, they might not find a
electron immediately.

Some of the holes can reach the depletion region towards our NMOS, and
be swept across the junction.

\subsubsection{Forward biased NMOS source or drain
junction}\label{forward-biased-nmos-source-or-drain-junction}

The increase in hole concentration underneath the NMOS can forward bias
the PN diode between source (or drain) and bulk. If this happens, then
we get electron injection into bulk. Some of those electrons can reach
the N-well depletion region, and be swept across (3).

\subsubsection{Positive-feedback}\label{positive-feedback}

Now we have a condition where the process accelerates, and locks-up.
Once turned on, this circuit will not turn off until the supply is low.

This is a phenomenon called latch-up. Similar to ESD circuits, latch-up
can short the supply to ground, and make things burn.

That is why, when we have digital logic, we need to be extra careful
close to the connection to the real world. Latch-up is bad.

We can prevent latch-up if we ensure that the electrons that start the
process never reach the N-wells. We can also prevent latch-up by
separating the NMOS and PMOS by guard rings (connections to ground, or
indeed supply), to serve as places where all these electrons and holes
can go.

Maybe it seems like a rare event for latch-up to happen, but trust me,
it's real, and it can happen in the strangest places. Similar to ESD,
it's a problem that can kill an IC, and make us pay another X million
dollars for a new tapeout, in addition to the layout work needed to fix
it.

Latch-up is why you will find the design rule check complaining if you
don't have enough substrate connections to ground, or N-well connections
to power close to your transistors.

Similar to the GGNMOS, this circuit, a
\href{https://en.wikipedia.org/wiki/Thyristor}{thyristor} can be a
useful circuit in ESD design. If we can trigger the thyristor when the
VDD shoots too high, then we can create a good ESD protection circuit.

See \href{https://www.sofics.com/features/low-leakage/}{low-leakage} ESD
for a few examples.

A model with the parasitic bipolars can be seen in Figure 16. The
resistors in the picture is to emulate what happens when there is a
current injected into the base of the NPN or PNP. I would recommend that
you think through the physics instead of using the parasitic bipolar
circuits. I've found the parasitic bipolar leads you down the wrong path
when you actually want to understand the physics of latch-up.


{
\centering
\aicfig{media/esd_scr_model_tikz.png}
}

\emph{Figure 16: Cross section of an inverter including the parasitic
bipolars }

You must \textbf{always handle ESD} on an IC

\begin{itemize}
\tightlist
\item
  Do everything yourself
\item
  Use libraries from foundry
\item
  Get help \href{http://www.sofics.com}{www.sofics.com}
\end{itemize}

\section{Summary}\label{summary}

The one-page version of this chapter:

\begin{itemize}
\tightlist
\item
  An ESD zap is a real event with a name: HBM is a charged person
  (kilovolts, amps through 1.5k), CDM is the chip discharging itself
  (nanoseconds, brutal)
\item
  Protection is a promise about every ordered pair of pads: each of the
  six permutations on a three-pad chip needs somewhere for the current
  to go
\item
  Two diodes per pin plus one rail clamp cover all permutations - adding
  a pin costs two diodes, not six paths
\item
  The grounded gate NMOS is off by every model you were taught; the
  parasitic lateral bipolar is what sinks the amperes
\item
  Thin gate oxide cannot take the leftover voltage, so real inputs add
  secondary protection behind a resistor
\item
  Latch-up is the same parasitic bipolars firing in normal operation:
  keep well taps close, keep injectors away from wells
\end{itemize}

\section{Would you like to know
more?}\label{would-you-like-to-know-more}

ESD (Electrostatic Discharge) Protection Design for Nanoelectronics in
CMOS Technology \citeproc{ref-ker06}{{[}3{]}}

Overview on Latch-Up Prevention in CMOS Integrated Circuits by Circuit
Solutions \citeproc{ref-ker23}{{[}4{]}}

Overview on ESD Protection Designs of Low-Parasitic Capacitance for RF
ICs in CMOS Technologies \citeproc{ref-ker11}{{[}5{]}}

\phantomsection\label{refs}
\begin{CSLReferences}{0}{0}
\bibitem[\citeproctext]{ref-huang21}
\CSLLeftMargin{{[}1{]} }%
\CSLRightInline{Z. Huang, Z. Tang, X.-P. Yu, Z. Shi, L. Lin, and N. N.
Tan, {``A BJT-based CMOS temperature sensor with duty-cycle-modulated
output and ±0.5°c (3\(\sigma\)) inaccuracy from \(-\)40 °c to 125 °c,''}
\emph{IEEE Transactions on Circuits and Systems II: Express Briefs},
vol. 68, no. 8, pp. 2780--2784, 2021, doi:
\href{https://doi.org/10.1109/TCSII.2021.3068283}{10.1109/TCSII.2021.3068283}.
}

\bibitem[\citeproctext]{ref-ker09}
\CSLLeftMargin{{[}2{]} }%
\CSLRightInline{M.-D. Ker, W.-Y. Chen, W.-T. Shieh, and I.-J. Wei,
{``New ballasting layout schemes to improve ESD robustness of i/o
buffers in fully silicided CMOS process,''} \emph{IEEE Transactions on
Electron Devices}, vol. 56, no. 12, pp. 3149--3159, 2009, doi:
\href{https://doi.org/10.1109/TED.2009.2031003}{10.1109/TED.2009.2031003}.
}

\bibitem[\citeproctext]{ref-ker06}
\CSLLeftMargin{{[}3{]} }%
\CSLRightInline{M.-D. Ker, {``ESD (electrostatic discharge) protection
design for nanoelectronics in CMOS technology,''} in \emph{2006 advanced
signal processing, circuit and system design techniques for
communications}, 2006, pp. 217--279, doi:
\href{https://doi.org/10.1109/ASPCAS.2006.251127}{10.1109/ASPCAS.2006.251127}.
}

\bibitem[\citeproctext]{ref-ker23}
\CSLLeftMargin{{[}4{]} }%
\CSLRightInline{M.-D. Ker and Z.-H. Jiang, {``Overview on latch-up
prevention in CMOS integrated circuits by circuit solutions,''}
\emph{IEEE Journal of the Electron Devices Society}, vol. 11, pp.
141--152, 2023, doi:
\href{https://doi.org/10.1109/JEDS.2022.3231822}{10.1109/JEDS.2022.3231822}.
}

\bibitem[\citeproctext]{ref-ker11}
\CSLLeftMargin{{[}5{]} }%
\CSLRightInline{M.-D. Ker, C.-Y. Lin, and Y.-W. Hsiao, {``Overview on
ESD protection designs of low-parasitic capacitance for RF ICs in CMOS
technologies,''} \emph{IEEE Transactions on Device and Materials
Reliability}, vol. 11, no. 2, pp. 207--218, 2011, doi:
\href{https://doi.org/10.1109/TDMR.2011.2106129}{10.1109/TDMR.2011.2106129}.
}

\end{CSLReferences}
